\subsection{Contexto del sistema y actores externos}

La Figura~\ref{fig:diagrama_de_contexto} delimita Carer-Record como una unidad y muestra las entidades externas que interactúan con la solución.

\begin{figure}[H]
    \centering
    \resizebox{0.95\textwidth}{!}{%
    \begin{tikzpicture}[
        actor/.style={draw, rounded corners, align=center, minimum width=3.8cm, minimum height=1.4cm, fill=gray!12},
        system/.style={draw=blue!60!black, very thick, rounded corners, align=center, minimum width=4.2cm, minimum height=2cm, fill=blue!8},
        flow/.style={-{Latex[length=3mm]}, thick}
    ]
        \node[actor] (clinical) {Personal asistencial\\Médicos y enfermería};
        \node[system, right=4.2cm of clinical] (carer) {Carer-Record\\Sistema de trazabilidad asistencial};
        \node[actor, right=4.2cm of carer] (source) {FCE simulada\\FHIR R4 en Docker, datos sintéticos};

        \draw[flow] ([yshift=4mm]clinical.east) -- node[above, align=center]{Consulta, registra notas\\y revisa propuestas} ([yshift=4mm]carer.west);
        \draw[flow] ([yshift=-4mm]carer.west) -- node[below, align=center]{Contexto autorizado,\\bitácora y resúmenes} ([yshift=-4mm]clinical.east);
        \draw[flow] (source.west) -- node[above, align=center]{Consulta mediante\\pasarela autorizada} (carer.east);
    \end{tikzpicture}%
    }
    \caption{Contexto y fronteras externas de Carer-Record}
    \label{fig:diagrama_de_contexto}
\end{figure}

El personal asistencial consulta información autorizada, registra notas propias y revisa propuestas. La aplicación accede a fuentes hospitalarias mediante una pasarela autorizada, no directamente. La demostración utiliza la FCE simulada FHIR R4 en Docker con datos sintéticos. Carer-Record no escribe en la FCE.
